% UNIFIED 2026-05-13
% SOURCE MAP for §4.3 (Plan 2, Agent 11):
%   - Statement of scaled Fusco--Julin (Thm 4.3.1): CITE(FJ2014, Thm 1.1) + DRAFT
%     scaling; reads Manuscript/00_notation_register.md §1, §7, §10;
%     source_map.md C2, C3, E3.1.
%   - Borel centre selection (Lem 4.3.2): DRAFT, with measurable-selection
%     prelim Agent 2 (source_map.md A7) and bulk-centroid Borel regularity
%     from Plan 2/wave3-F-h1-center-bound.md §2.2 (LOCAL).
%   - Fraenkel bound (Lem 4.3.3): DRAFT, consumes FJ scaled form (C3) and
%     FMP/AFM/CL centre proximity (LOCAL, wave3-F-h1-center-bound.md §1
%     + Wave 2 Agent A Lemma 2.2).
%   - Normal-oscillation bound (Lem 4.3.4): DRAFT, divergence identity
%     Agent 4 §C4; FJ2014 Thm 1.1.
%   - Oriented radial trace lemma (Thm 4.3.5, the critical load-bearing
%     lemma E3.5): DRAFT, builds on Plan 2/WIP/WIP_WeightedMetricTrace.tex
%     Lem.\ 3.6 and the wave3-J-route-delta-repair.md §§3--7 slicing
%     decomposition; uses Agent 2 prelims A4 (divergence theorem for
%     finite-perimeter sets, CITE Maggi2012 Thm 16.3) and A5 (Lipschitz
%     truncation of x/|x|, DRAFT in §2.1).
%   - Vector field X(x)=||x-z|/rho-1|(x-z)/|x-z|: named precisely as in
%     the brief and in WIP_WeightedMetricTrace.tex (LOCAL).
%   - All constants C(n, R, rho_*) tracked per notation register §14.

\subsection{\S4.3.\ Fusco--Julin centre and the oriented radial trace lemma}
\label{sec:plan2-fj-radial}

This subsection is the centre-package and load-bearing radial-trace
estimate of Plan~2. Its conclusion --- Theorem~\ref{thm:radial-trace} ---
together with the metric first-variation theorem of \S4.2 controls the
metric derivative $|\dot F_\rho|_{\X}$ in the weighted action integral.

Throughout this subsection we fix $n\ge 2$, $R\ge 2$, $\rstar\in(0,1)$,
and $\kappa>0$. We work with an open set $\Om\subset B_R\subset\R^n$,
$|\Om|=\omega_n$, with variational torsion function $u=u_\Om$, volume-radius
parametrisation $E_\rho=\{u>t(\rho)\}$, $|E_\rho|=\omega_n\rho^n$ for
$\rho\in[\rstar,1]$, and $\mu(d\rho):=(-t'(\rho))\,d\rho$. We assume
\begin{equation}\label{eq:fj-smallness}
   \delta_T(\Om)\le\delta_0(n,R,\rstar),
\end{equation}
where $\delta_0$ is the smallness threshold fixed in
Theorem~\ref{thm:weighted-trace} (the constant matching the
Cicalese--Leonardi smallness threshold and the centroid--Fraenkel offset,
cf.\ \cite[Prop.~2.3]{CL2012} and \cite[Thm.~1.1]{FJ2014}).

\subsubsection{4.3.1.\ Scaled Fusco--Julin and centre selection}
\label{ssec:fj-scaled}

Recall the strong quantitative isoperimetric inequality of
Fusco--Julin~\cite[Thm.\ 1.1]{FJ2014}: there is a dimensional constant
$C_{\mathrm{FJ}}(n)>0$ such that for every set of finite perimeter
$E\subset\R^n$ with $|E|=|B_1|$ there exists a centre $z=z(E)\in\R^n$ with
\begin{equation}\label{eq:FJ-original}
   \Bigl(\frac{|E\Delta B_1(z)|}{|B_1|}\Bigr)^2
   +\beta(E,z)^2
   \le C_{\mathrm{FJ}}(n)\,\bigl(P(E)-P(B_1)\bigr),
\end{equation}
where the FJ oscillation index is
\begin{equation}\label{eq:FJ-osc-def}
   \beta(E,z)^2
   :=\inf_{|E'|=|E|}\;
   \int_{\partial^*E}\bigl|\nu_E(x)-(\nu_{E'})_{\mathrm{rad}}(x;z)\bigr|^2\,d\mathcal H^{n-1}(x),
\end{equation}
which by the standard divergence identity (item~C4 of the
preliminaries; see~\eqref{eq:beta-div} below) is comparable to
$\int_{\partial^*E}|\nu_E-(x-z)/|x-z||^2\,d\mathcal H^{n-1}$.

\begin{theorem}[Scaled Fusco--Julin on $\{|E|=\omega_n\rho^n\}$]
\label{thm:plan2-FJ-scaled}\label{prop:fj-scaled}
For every $\rho\in[\rstar,1]$ and every $E\subset\R^n$ of finite perimeter
with $|E|=\omega_n\rho^n$ there exists $z=z(E,\rho)\in\R^n$ such that
\begin{equation}\label{eq:plan2-FJ-scaled}
   |E\Delta B_\rho(z)|^2
   +\rho^{n+1}\int_{\partial^*E}\!\!\Bigl|\nu_E(x)-\frac{x-z}{|x-z|}\Bigr|^2 d\mathcal H^{n-1}(x)
   \le C_1(n,\rstar)\,\rho^{n+1}\bigl(P(E)-P(B_\rho)\bigr),
\end{equation}
where $C_1(n,\rstar)=4\,C_{\mathrm{FJ}}(n)\cdot\max(1,\rstar^{-(n+1)})$. The
constant depends only on $n$ and $\rstar$.
\end{theorem}

\begin{proof}
Set $\widetilde E:=\rho^{-1}E$, so $|\widetilde E|=\omega_n=|B_1|$ and
$\widetilde E$ has finite perimeter with
$P(\widetilde E)=\rho^{-(n-1)}P(E)$. Apply~\eqref{eq:FJ-original} to
$\widetilde E$: there exists $\widetilde z$ with
\begin{equation}\label{eq:FJ-applied}
   \Bigl(\frac{|\widetilde E\Delta B_1(\widetilde z)|}{\omega_n}\Bigr)^2
   +\beta(\widetilde E,\widetilde z)^2
   \le C_{\mathrm{FJ}}(n)\bigl(P(\widetilde E)-P(B_1)\bigr).
\end{equation}
Set $z:=\rho\widetilde z$. Translating $\widetilde E$ to $E$ scales
volumes by $\rho^n$ and perimeters by $\rho^{n-1}$, and using the
divergence identity~\eqref{eq:beta-div} below to rewrite
$\beta(\widetilde E,\widetilde z)^2$ in terms of the
normal-oscillation integral, the scaled bounds
$|E\Delta B_\rho(z)|=\rho^n|\widetilde E\Delta B_1(\widetilde z)|$
and $\int_{\partial^*E}|\nu_E-(x-z)/|x-z||^2\,d\mathcal H^{n-1}
=\rho^{n-1}\int_{\partial^*\widetilde E}|\nu_{\widetilde E}-(y-\widetilde z)/|y-\widetilde z||^2\,d\mathcal H^{n-1}$
yield~\eqref{eq:plan2-FJ-scaled} after multiplying through by $\rho^{2n}$
and using $\rstar\le\rho\le1$ to absorb the $\rho$-powers into the
constant $C_1$.
\end{proof}

\begin{remark}[Comparison constants]\label{rmk:FJ-equiv}
The divergence identity (item~C4 of the preliminaries)
\begin{equation}\label{eq:beta-div}
   2\,\beta(E,z)^2
   = \int_{\partial^*E}\Bigl|\nu_E(x)-\frac{x-z}{|x-z|}\Bigr|^2 d\mathcal H^{n-1}(x)
   = 2(P(E)-P(B_\rho))+2\Pi(E,z),
\end{equation}
holds whenever $|E|=|B_\rho|$, $z\in\R^n$ with the radial field
$x\mapsto(x-z)/|x-z|$ locally Borel and the truncation argument
of~\S2.1 (item~A5) applicable. The identity is proved by testing
$\mathrm{div}((x-z)/|x-z|)=(n-1)|x-z|^{-1}$ on $E$ via the finite-perimeter
divergence theorem and truncating away from~$z$;
$\Pi(E,z):=(n-1)\int (\mathbf1_{B_\rho(z)}-\mathbf1_E)|x-z|^{-1}\,dx$ is the
signed radial moment of the notation register \S13. We use only the
inequality
$\int|\nu_E-(x-z)/|x-z||^2\,d\mathcal H^{n-1}\le 2C_1(n,\rstar)(P(E)-P(B_\rho))/\rho^{n+1}$
that follows from~\eqref{eq:plan2-FJ-scaled}.
\end{remark}

We now choose a single Borel centre field along the torsion foliation.
The canonical choice (see~\S\ref{sec:notation}, and the discussion of the
Fusco--Julin centre in~\cite[Thm.~1.1]{FJ2014})
is the bulk centroid
\begin{equation}\label{eq:centroid-def}
   \zbar
   = \bar z_\rho
   :=\frac{1}{|E_\rho|}\int_{E_\rho}x\,dx
   \in\R^n,\qquad\rho\in[\rstar,1].
\end{equation}

\begin{lemma}[Borel measurability and CL containment of the centroid]
\label{lem:centroid-borel}
Under~\eqref{eq:fj-smallness} the map $\rho\mapsto\zbar$ is Lipschitz on
$[\rstar,1]$ with constant $L_{\mathrm{cent}}(n,R,\rstar)\le 2nR\rstar^{-n}$;
in particular it is Borel. Moreover for every $\rho\in[\rstar,1]$,
\begin{equation}\label{eq:zbar-CL}
   B(\zbar,\rho/4)\subset E_\rho\subset B(\zbar,3\rho),
\end{equation}
so $\zbar\in E_\rho$ and the radial field $e_r(x):=(x-\zbar)/|x-\zbar|$ is
$\mathcal H^{n-1}$-a.e.\ well-defined and Borel on $\partial^*E_\rho$, with
$\rstar/4\le|x-\zbar|\le 3R$ on $\partial^*E_\rho$ for a.e.\ $\rho$.
\end{lemma}

\begin{proof}
For $\rstar\le\rho\le\rho'\le1$, the nesting of $\{E_\rho\}$ gives
$|E_\rho\Delta E_{\rho'}|=|E_{\rho'}|-|E_\rho|\le n\omega_n|\rho-\rho'|$.
Hence
$|\int_{E_\rho}x_i\,dx-\int_{E_{\rho'}}x_i\,dx|\le R\cdot n\omega_n|\rho-\rho'|$;
combined with $|E_\rho|\ge\omega_n\rstar^n$ this yields the Lipschitz
constant $L_{\mathrm{cent}}=2nR/\rstar^n$, which is Borel (in fact
continuous). The containment~\eqref{eq:zbar-CL} follows from the
Cicalese--Leonardi inner-ball containment (\cite[Thm.\ 1.1]{CL2012})
applied at the Fraenkel centre
$z^{\mathrm{Fr}}_\rho$, together with the centroid--Fraenkel proximity
$|\zbar-z^{\mathrm{Fr}}_\rho|\le C(n,\rstar)\,\Fr(E_\rho)\le
C'\sqrt{\delta_T(\Om)}\le\rho/4$, valid for $\delta_0$ small
(\cite[Prop.\ 2.3]{CL2012},
\cite[Lem.\ 2.6]{AFM2013}, \cite[Thm.~1.1]{FJ2014}).
Borel measurability of the normal field $\nu_{E_\rho}$ on $\partial^*E_\rho$
(\cite[Thm.\ 15.9]{Maggi2012}) then yields the Borel measurability of
$e_r$ on $\partial^*E_\rho$ for $\zbar\in E_\rho$.
\end{proof}

\begin{remark}[Auxiliary FJ-oscillation centre and Borel selection]
\label{rmk:centre-selection}
The centre $z(E_\rho,\rho)$ produced by Theorem~\ref{thm:plan2-FJ-scaled} is
\emph{a} valid almost-minimiser; Lemma~\ref{lem:centroid-borel}'s centroid
$\zbar$ also satisfies~\eqref{eq:plan2-FJ-scaled} up to enlarging $C_1$ by the
proximity offset $|\zbar-z^{\mathrm{Fr}}_\rho|\le C(n,\rstar)\sqrt{\delta_T}$,
absorbed by the smallness threshold $\delta_0$. The Borel selection
of \emph{any} valid centre on a Borel-measurable family of finite-perimeter
sets follows from the measurable-selection statement A7 of the
preliminaries (Federer~\cite[\S2.13]{Federer1969}, in the form recorded
in~\cite[\S5.3]{AFP2000}). Throughout the rest of this subsection the
centre is fixed to be $z_\rho:=\zbar$.
\end{remark}

\subsubsection{4.3.2.\ Fraenkel and normal-oscillation bounds at the centroid}
\label{ssec:fj-bounds}

Define the (good) isoperimetric set
\begin{equation}\label{eq:GI-def}
   G_I:=\{\rho\in[\rstar,1]:D_I(t(\rho))\le\eta_I\},\qquad
   \eta_I:=\eta_I(n,R,\rstar)\in(0,1],
\end{equation}
with $\eta_I$ small enough that on $G_I$, $P(E_\rho)\le 2P(B_\rho)$
(possible since $D_I=(P(E)-P(B_\rho))(P(E)+P(B_\rho))$ and
$P(B_\rho)\ge n\omega_n\rstar^{n-1}$). We collect
\eqref{eq:plan2-FJ-scaled} at $z_\rho=\zbar$.

\begin{lemma}[Fraenkel bound at the centroid]\label{lem:fraenkel-centroid}
There exists $C_2=C_2(n,R,\rstar)>0$ such that for every $\rho\in G_I$,
\begin{equation}\label{eq:fraenkel-centroid}
   |E_\rho\Delta B_\rho(\zbar)|^2
   \le C_2(n,R,\rstar)\,D_I(t(\rho)).
\end{equation}
\end{lemma}

\begin{proof}
From~\eqref{eq:plan2-FJ-scaled} applied with $z=\zbar$ (valid by
Remark~\ref{rmk:centre-selection}, after absorbing the
$|\zbar-z^{\mathrm{Fr}}_\rho|^2$ correction into the constant via
$|B_\rho(a)\Delta B_\rho(b)|\le C(n)\rho^{n-1}|a-b|$ and the FMP estimate
$|E_\rho\Delta B_\rho(z^{\mathrm{Fr}}_\rho)|^2\le c_1|B_\rho|^2(P(E_\rho)-P(B_\rho))$),
\begin{equation}\label{eq:fraenkel-centroid-pre}
   |E_\rho\Delta B_\rho(\zbar)|^2
   \le 2C_1(n,\rstar)\rho^{n+1}\bigl(P(E_\rho)-P(B_\rho)\bigr).
\end{equation}
Use the algebraic identity
$D_I(t(\rho))=(P(E_\rho)-P(B_\rho))(P(E_\rho)+P(B_\rho))\ge n\omega_n\rstar^{n-1}\cdot(P(E_\rho)-P(B_\rho))$
to convert $(P(E_\rho)-P(B_\rho))$ into $D_I/(n\omega_n\rstar^{n-1})$;
this and $\rho\le1$ yield~\eqref{eq:fraenkel-centroid} with
$C_2=2C_1(n,\rstar)/(n\omega_n\rstar^{n-1})$, polynomial in $n,R,1/\rstar$.
\end{proof}

\begin{lemma}[Normal-oscillation bound at the centroid]
\label{lem:normal-osc-centroid}
There exists $C_3=C_3(n,R,\rstar)>0$ such that for every $\rho\in G_I$,
\begin{equation}\label{eq:normal-osc-centroid}
   \int_{\partial^*E_\rho}\!\!\Bigl|\nu_\rho(x)-\frac{x-\zbar}{|x-\zbar|}\Bigr|^2\,
      d\mathcal H^{n-1}(x)
   \le C_3(n,R,\rstar)\,D_I(t(\rho)).
\end{equation}
\end{lemma}

\begin{proof}
Combining~\eqref{eq:plan2-FJ-scaled} at $z=\zbar$ (with the same absorption as
in the previous proof) and the lower bound $\rho^{n+1}\ge\rstar^{n+1}$
gives
$\int|\nu_\rho-(x-\zbar)/|x-\zbar||^2\,d\mathcal H^{n-1}\le
C_1(n,\rstar)(P(E_\rho)-P(B_\rho))/\rstar^{n+1}$; converting $P-P_\rho$ to
$D_I$ as in the previous proof yields~\eqref{eq:normal-osc-centroid}
with $C_3=C_1(n,\rstar)/(\rstar^{n+1}\cdot n\omega_n\rstar^{n-1})$.
\end{proof}

\begin{remark}[{$\rho\in[\rstar,1]\setminus G_I$}]
Outside $G_I$, $D_I(t(\rho))>\eta_I$, and~\eqref{eq:fraenkel-centroid},
\eqref{eq:normal-osc-centroid} hold trivially after enlarging
$C_2,C_3$ by a multiple of $\eta_I^{-1}$, since both left-hand sides are
bounded by $|B_R|^2$ and $4P(B_R)\le 4n\omega_n R^{n-1}$ respectively.
In the integrated argument of Theorem~\ref{thm:weighted-trace} the
$\mu$-measure of $\{\rho:\rho\notin G_I\}$ is controlled by Markov from
the deficit identity, so the trivial bound is permissible.
\end{remark}

\subsubsection{4.3.3.\ The oriented radial trace lemma}
\label{ssec:radial-trace}

This is the load-bearing lemma of Plan~2's GMT route. We prove the
pointwise inequality, with the named vector field
\begin{equation}\label{eq:X-def}
   X(x)
   :=\Bigl|\frac{|x-\zbar|}{\rho}-1\Bigr|\cdot\frac{x-\zbar}{|x-\zbar|},
   \qquad x\in\R^n\setminus\{\zbar\},
\end{equation}
of the brief, via the finite-perimeter divergence theorem (item~A4 of
the preliminaries: \cite[Thm.\ 16.3]{Maggi2012}) and the truncation
argument of item~A5. We do \emph{not} use ray-slicing in the proof of
this lemma; a slicing-based variant is sketched in
Remark~\ref{rmk:slicing-route} below.

\begin{theorem}[Oriented $L^1$ radial trace at the centroid]
\label{thm:radial-trace}\label{thm:fj-radial-divergence}\label{lem:radial-trace-master}
There is $C_4=C_4(n,R,\rstar)>0$ such that for every
$\rho\in[\rstar,1]$ satisfying~\eqref{eq:fj-smallness}, with the centroid
$z_\rho=\zbar$ of~\eqref{eq:centroid-def},
\begin{equation}\label{eq:radial-trace}
   T_\rho
   :=\int_{\partial^*E_\rho}\!\!\Bigl|\frac{|x-\zbar|}{\rho}-1\Bigr|\,
       d\mathcal H^{n-1}(x)
   \le
   C_4\,\Bigl(\,|E_\rho\Delta B_\rho(\zbar)|
   +\int_{\partial^*E_\rho}\!\!\Bigl|\nu_\rho-\frac{x-\zbar}{|x-\zbar|}\Bigr|^2\,
       d\mathcal H^{n-1}\Bigr).
\end{equation}
The constant has the explicit form $C_4(n,R,\rstar)=\max\{n/\rstar,3R/(2\rstar)\}$
(see~\eqref{eq:T-trace-explicit} below).
\end{theorem}

\begin{proof}
\emph{Step 1: setup.} By Lemma~\ref{lem:centroid-borel},
$\zbar\in B(\zbar,\rho/4)\subset E_\rho\subset B(\zbar,3\rho)\subset B_R$,
so $|x-\zbar|\in[\rstar/4,3R]$ for $\mathcal H^{n-1}$-a.e.
$x\in\partial^*E_\rho$. Translate so that $\zbar=0$; the radial unit
vector is $e_r(x)=x/|x|$ on $\R^n\setminus\{0\}$ and the
weight is $w(r):=|r/\rho-1|$, a Lipschitz function of $r\in[0,\infty)$
with $w(\rho)=0$, $|w(r)|\le \max(1,(3R)/\rho-1)\le 3R/\rstar$ on
$\partial^*E_\rho$, and the
vector field $X(x)=w(|x|)e_r(x)$ of~\eqref{eq:X-def}. Note $X$ is locally
Lipschitz on $\R^n\setminus\{0\}$ and bounded on $B_R$ (with
$\|X\|_\infty\le 3R/\rstar$), but is discontinuous at the origin.

\emph{Step 2: pointwise inequality.}
On $\partial^*E_\rho$, $|e_r|=|\nu_\rho|=1$, so writing
$1=e_r\cdot\nu_\rho+(1-e_r\cdot\nu_\rho)$ and using $w\ge0$,
\begin{equation}\label{eq:T-split}
   T_\rho
   =\int_{\partial^*E_\rho}w(|x|)\,d\mathcal H^{n-1}
   =\int_{\partial^*E_\rho}X\cdot\nu_\rho\,d\mathcal H^{n-1}
    +\int_{\partial^*E_\rho}w(|x|)\bigl(1-e_r\cdot\nu_\rho\bigr)\,d\mathcal H^{n-1}.
\end{equation}
The second integrand is non-negative because $e_r\cdot\nu_\rho\le|e_r||\nu_\rho|=1$,
and the elementary identity
$1-e_r\cdot\nu_\rho=\tfrac12|e_r-\nu_\rho|^2$ together with
$w\le3R/\rstar$ gives
\begin{equation}\label{eq:T-tan-bound}
   \int_{\partial^*E_\rho}w(|x|)(1-e_r\cdot\nu_\rho)\,d\mathcal H^{n-1}
   \le\frac{3R}{2\rstar}
   \int_{\partial^*E_\rho}|\nu_\rho-e_r|^2\,d\mathcal H^{n-1}.
\end{equation}

\emph{Step 3: truncated divergence theorem for the radial term.}
We bound $\int_{\partial^*E_\rho}X\cdot\nu_\rho\,d\mathcal H^{n-1}$ by the
finite-perimeter divergence theorem applied to a Lipschitz truncation of
$X$. Fix $\eps\in(0,\rstar/8)$ and a smooth non-decreasing
$\eta_\eps:[0,\infty)\to[0,1]$ with $\eta_\eps(r)=0$ for $r\le\eps/2$,
$\eta_\eps(r)=1$ for $r\ge\eps$, $|\eta_\eps'|\le4/\eps$. Set
\begin{equation}\label{eq:Xeps-def}
   X_\eps(x):=\eta_\eps(|x|)\,w(|x|)\,e_r(x),\qquad x\in\R^n.
\end{equation}
Then $X_\eps\in C^{0,1}(\R^n;\R^n)\cap C^\infty(\R^n\setminus\{0\};\R^n)$,
$X_\eps\equiv X$ outside $\{|x|<\eps\}$, and
$\|X_\eps\|_\infty\le 3R/\rstar$. The classical divergence
$\mathrm{div}\,X_\eps$ is locally bounded on $\R^n$ and equals
\begin{equation}\label{eq:divXeps}
   \mathrm{div}\,X_\eps(x)
   = \eta_\eps(|x|)\Bigl[w'(|x|)+\frac{n-1}{|x|}w(|x|)\Bigr]
     +\eta_\eps'(|x|)\,w(|x|)\qquad(x\ne0),
\end{equation}
where $w'(r)=r/(\rho|r|)\cdot\mathrm{sgn}(r/\rho-1)=\rho^{-1}\mathrm{sgn}(r/\rho-1)$
for $r\ne\rho$ (Lipschitz regularity of $w$, derivative a.e.).
At the origin, $X_\eps\equiv 0$ on $\{|x|<\eps/2\}$, so $X_\eps$ is genuinely
Lipschitz on $\R^n$.

By the divergence theorem for the finite-perimeter set $E_\rho$ applied to
the Lipschitz vector field $X_\eps$ (item~A4 of the preliminaries;
\cite[Thm.\ 16.3]{Maggi2012}, applied with the standard density argument:
$X_\eps$ is mollifiable since it is Lipschitz with compact support after
multiplication by a cutoff at radius $3R$),
\begin{equation}\label{eq:div-eps}
   \int_{\partial^*E_\rho}X_\eps\cdot\nu_\rho\,d\mathcal H^{n-1}
   =\int_{E_\rho}\mathrm{div}\,X_\eps\,dx.
\end{equation}
The same identity applied to the open ball $B_\rho=B_\rho(0)$ (a smooth
finite-perimeter set, with $\nu_{B_\rho}(y)=y/|y|$ and $w(\rho)=0$) yields
\begin{equation}\label{eq:div-eps-ball}
   0
   =\int_{\partial^*B_\rho}\!\!\!\!w(|y|)\,e_r\cdot\nu_{B_\rho}\,d\mathcal H^{n-1}
   =\int_{\partial^*B_\rho}\!\!\!\!X_\eps\cdot\nu_{B_\rho}\,d\mathcal H^{n-1}
   =\int_{B_\rho}\!\!\mathrm{div}\,X_\eps\,dx,
\end{equation}
where the first equality holds because $\eta_\eps(r)=1$ for $r$ near $\rho$
(since $\eps<\rstar/8\le\rho/8$) and $w(\rho)=0$. Subtracting,
\begin{equation}\label{eq:div-eps-diff}
   \int_{\partial^*E_\rho}X_\eps\cdot\nu_\rho\,d\mathcal H^{n-1}
   =\int_{E_\rho\setminus B_\rho}\!\!\!\!\mathrm{div}\,X_\eps\,dx
    -\int_{B_\rho\setminus E_\rho}\!\!\!\!\mathrm{div}\,X_\eps\,dx.
\end{equation}

\emph{Step 4: limit $\eps\downarrow0$.} On $\R^n\setminus\{0\}$,
$X_\eps\to X$ pointwise and $\mathrm{div}\,X_\eps\to\mathrm{div}\,X$
pointwise, where
\begin{equation}\label{eq:divX}
   \mathrm{div}\,X(x)=
   \rho^{-1}\mathrm{sgn}(|x|/\rho-1)+\frac{n-1}{|x|}\Bigl|\frac{|x|}{\rho}-1\Bigr|,
   \qquad x\in\R^n\setminus(\{0\}\cup\partial B_\rho).
\end{equation}
On both $E_\rho\setminus B_\rho$ and $B_\rho\setminus E_\rho$, the
domain of integration is contained in $\{|x|\ge\rstar/4\}$ (by
Lemma~\ref{lem:centroid-borel} for the first set and the trivial bound
$\rho\ge\rstar$ for the second after restricting to $|x|\ge\rstar/4$;
the set $\{|x|<\rstar/4\}\cap(E_\rho\cup B_\rho)\subset B_\rho$ is
\emph{not} in either symmetric difference part by~\eqref{eq:zbar-CL}).
On $\{|x|\ge\rstar/4\}$, the singular factor
$(n-1)/|x|\le 4(n-1)/\rstar$ and the cutoff factor $\eta_\eps(|x|)\equiv1$
for $\eps<\rstar/8$; therefore $X_\eps\equiv X$ on
$\{|x|\ge\rstar/4\}$ for all $\eps$ small enough, and
\eqref{eq:div-eps-diff} is in fact independent of $\eps\in(0,\rstar/8)$
and equals
\begin{equation}\label{eq:div-final}
   \int_{\partial^*E_\rho}X\cdot\nu_\rho\,d\mathcal H^{n-1}
   =\int_{E_\rho\setminus B_\rho}\!\!\!\mathrm{div}\,X\,dx
    -\int_{B_\rho\setminus E_\rho}\!\!\!\mathrm{div}\,X\,dx.
\end{equation}
(The left-hand side is $\int X\cdot\nu_\rho\,d\mathcal H^{n-1}$, the limit
of the same integral with $X_\eps$ in place of $X$, by the same argument
applied to the boundary trace.) This is the truncation step (A5 of the
preliminaries): the locally integrable singularity of $X$ at $\zbar=0$
is harmless because $\zbar\in\mathrm{int}\,E_\rho$ and the domains
$E_\rho\setminus B_\rho$, $B_\rho\setminus E_\rho$ are
\emph{disjoint from a fixed neighbourhood} of $\zbar$.

\emph{Step 5: pointwise bound on $\mathrm{div}\,X$.}
For $r:=|x|>\rho$ (i.e.\ $x\in E_\rho\setminus B_\rho$):
\begin{equation}\label{eq:divX-out}
   \mathrm{div}\,X(x)
   =\rho^{-1}+\frac{n-1}{r}\Bigl(\frac{r}{\rho}-1\Bigr)
   =\frac{n}{\rho}-\frac{n-1}{r}
   \le\frac{n}{\rho}\le\frac{n}{\rstar}.
\end{equation}
For $r<\rho$ (i.e.\ $x\in B_\rho\setminus E_\rho$, $r\ge\rstar/4$):
\begin{equation}\label{eq:divX-in}
   \mathrm{div}\,X(x)
   =-\rho^{-1}+\frac{n-1}{r}\Bigl(1-\frac{r}{\rho}\Bigr)
   =\frac{n-1}{r}-\frac{n}{\rho},
\end{equation}
so
$(-\mathrm{div}\,X)_+(x)\le n/\rho\le n/\rstar$. Combining with
\eqref{eq:div-final},
\begin{equation}\label{eq:T-rad-bound}
   \int_{\partial^*E_\rho}X\cdot\nu_\rho\,d\mathcal H^{n-1}
   \le\frac{n}{\rstar}\,|E_\rho\Delta B_\rho(\zbar)|.
\end{equation}

\emph{Step 6: assembly.}
Combining~\eqref{eq:T-split}, \eqref{eq:T-tan-bound},
and~\eqref{eq:T-rad-bound},
\begin{equation}\label{eq:T-trace-explicit}
   T_\rho
   \le\frac{n}{\rstar}\,|E_\rho\Delta B_\rho(\zbar)|
   +\frac{3R}{2\rstar}\,
     \int_{\partial^*E_\rho}|\nu_\rho-e_r|^2\,d\mathcal H^{n-1}.
\end{equation}
This is~\eqref{eq:radial-trace} with
$C_4(n,R,\rstar)=\max\{n/\rstar,\,3R/(2\rstar)\}$, polynomial in $n,R,1/\rstar$.
\end{proof}

\begin{remark}[Why the truncation is legitimate]\label{rmk:truncation}
The truncation factor $\eta_\eps$ is needed only to ensure that the
classical divergence theorem applies to a Lipschitz field on all of
$\R^n$; the field $X$ is bounded on $B_R$, but its classical derivative
$\mathrm{div}\,X$ has a singularity of type $|x|^{-1}$ at $\zbar$. Two
properties make the truncation harmless. First,
$\zbar\in\mathrm{int}\,E_\rho$ (Lemma~\ref{lem:centroid-borel}), so the
small ball $\{|x-\zbar|<\rstar/4\}$ lies entirely in $E_\rho\cap B_\rho$
and contributes \emph{equally} to both $E_\rho$ and $B_\rho$ in
\eqref{eq:div-eps-diff}, cancelling out. Second, on the
symmetric-difference set $\{x:\mathbf1_{E_\rho}(x)\ne\mathbf1_{B_\rho}(x)\}$
the radius $|x-\zbar|\ge\rstar/4$ (since this set lies in
$E_\rho\Delta B_\rho\subset B_R\setminus B(\zbar,\rstar/4)$ once
$\delta_0$ is small), so $X=X_\eps$ on this set for all $\eps<\rstar/8$
and the limit~\eqref{eq:div-final} is exact, not a passage. This is the
key estimate of the truncation step: the singularity of $X$ at $\zbar$
contributes nothing because it is interior to both $E_\rho$ and $B_\rho$.
\end{remark}

\begin{remark}[Equivalent slicing route]\label{rmk:slicing-route}
An alternative proof of~\eqref{eq:radial-trace} by a spherical slicing
decomposition $T_\rho=T_\rho^{\mathrm{rad,slic}}+T_\rho^{\mathrm{tan}}$
can be carried out via the radial slicing formula
\cite[Thm.\ 18.11]{Maggi2012} and a good-ray/bad-ray multiplicity
decomposition. That route yields the same bound at the integrated level
and trades the divergence-theorem step here for the Cicalese--Leonardi
annular containment. Both routes produce the constant $C_4(n,R,\rstar)$
in the same polynomial form. The divergence-theorem route is preferred
here because it isolates the named vector field~\eqref{eq:X-def} and
avoids reference to angular multiplicity.
\end{remark}

\subsubsection{4.3.4.\ Consequences for the metric first-variation chain}
\label{ssec:fj-output}

We collect the levelwise outputs that feed into the integrated action
bound of Theorem~\ref{thm:weighted-trace} (\S4.4).

\begin{corollary}[Pointwise radial trace from $D_I$]\label{cor:T-pointwise}
There exists $C_5=C_5(n,R,\rstar)$ such that for every $\rho\in G_I$,
\begin{equation}\label{eq:T-pointwise}
   T_\rho^2
   \le C_5(n,R,\rstar)\,D_I(t(\rho)).
\end{equation}
\end{corollary}

\begin{proof}
By Theorem~\ref{thm:radial-trace},
$T_\rho^2\le 2C_4^2(|E_\rho\Delta B_\rho(\zbar)|^2
+(\int|\nu_\rho-e_r|^2\,d\mathcal H^{n-1})^2)$. The first term is
$\le 2C_4^2\cdot C_2\,D_I$ by Lemma~\ref{lem:fraenkel-centroid}. For the
second, use the trivial estimate
$(\int|\nu_\rho-e_r|^2\,d\mathcal H^{n-1})^2\le 4P(B_R)\cdot\int|\nu_\rho-e_r|^2\,d\mathcal H^{n-1}$
on $G_I$ (since the integrand is bounded by~$4$ pointwise and the
perimeter $P(E_\rho)\le 2P(B_\rho)\le 2n\omega_n R^{n-1}$ on $G_I$),
together with Lemma~\ref{lem:normal-osc-centroid}, to get
$\le 8n\omega_n R^{n-1}C_3\,D_I$. Combining,
$C_5=2C_4^2(C_2+8n\omega_n R^{n-1}C_3)$, polynomial in $n,R,1/\rstar$.
\end{proof}

\begin{corollary}[Good-level homothetic trace at the centroid]
\label{cor:H-trace-good}
With $\bar V_\rho:=P(B_\rho)/P(E_\rho)$ and
$H_{z,\rho}(x)=(x-z)\cdot\nu_\rho(x)/\rho$ on $\partial^*E_\rho$, there
exists $C_6=C_6(n,R,\rstar)$ such that for every $\rho\in G_I$,
\begin{equation}\label{eq:H-trace-good}
   \Bigl(\int_{\partial^*E_\rho}|H_{\zbar,\rho}-\bar V_\rho|\,
       d\mathcal H^{n-1}\Bigr)^2
   \le C_6(n,R,\rstar)\,D_I(t(\rho)).
\end{equation}
\end{corollary}

\begin{proof}
Write $e_r(x)=(x-\zbar)/|x-\zbar|$. Then
$H_{\zbar,\rho}(x)=(|x-\zbar|/\rho)\,e_r\cdot\nu_\rho$ and
\[
   |H_{\zbar,\rho}-1|
   \le\Bigl|\frac{|x-\zbar|}{\rho}-1\Bigr|+|e_r\cdot\nu_\rho-1|.
\]
The second term equals $\tfrac12|e_r-\nu_\rho|^2$. Integrating,
\[
   \int_{\partial^*E_\rho}|H_{\zbar,\rho}-1|\,d\mathcal H^{n-1}
   \le T_\rho+\frac12\int_{\partial^*E_\rho}|\nu_\rho-e_r|^2\,d\mathcal H^{n-1}.
\]
Squaring and using $(a+b)^2\le2a^2+2b^2$ together with
Corollary~\ref{cor:T-pointwise} and Lemma~\ref{lem:normal-osc-centroid}
gives
$(\int|H_{\zbar,\rho}-1|)^2\le C\,D_I$ on $G_I$. The replacement of
the centring constant $1$ by $\bar V_\rho=1-(P(E_\rho)-P(B_\rho))/P(E_\rho)$
costs an additive term
\[
   |1-\bar V_\rho|\cdot P(E_\rho)=P(E_\rho)-P(B_\rho)
   \le D_I(t(\rho))/(n\omega_n\rstar^{n-1})\le D_I^{1/2}/(n\omega_n\rstar^{n-1})
\]
on $G_I$ (using $D_I\le\eta_I\le1$ on $G_I$). Squaring and absorbing
this term yields~\eqref{eq:H-trace-good} with
$C_6=2(C_5+C_3)+2/(n\omega_n\rstar^{n-1})^2$, polynomial in $n,R,1/\rstar$.
\end{proof}

\begin{remark}[Hypothesis tracking]\label{rmk:hypothesis-tracking}
Every estimate above is stated for open $\Om\subset B_R$ with
$|\Om|=\omega_n$ and $\delta_T(\Om)\le\delta_0(n,R,\rstar)$ small enough
to ensure Cicalese--Leonardi smallness for every level
$\rho\in[\rstar,1]$ via \cite[Thm.~1.1]{CL2012}. No graph
parametrisation, no minimiser selection, and no smoothness of
$\partial^*E_\rho$ are used. All boundary integrals are on the De Giorgi
reduced boundary. The Borel selection of the centroid is unconditional
(Lemma~\ref{lem:centroid-borel}); the FJ centre selection of
Remark~\ref{rmk:centre-selection} reduces to it after the
$O(\sqrt{\delta_T})$ proximity absorption recorded in~\cite[Thm.~1.1]{FJ2014}.
Outside the smallness regime, the FMP fallback
$\Fr(\Om)^2\le C(n)\delta_T^{1/2}$ together with $\delta_T\ge c(n)>0$
gives all four estimates trivially.
\end{remark}

\paragraph*{Constants used in this subsection.} For the convenience of
\S4.4 we collect:
\begin{itemize}
\item $C_1(n,\rstar)$ in \eqref{eq:plan2-FJ-scaled}: scaled FJ constant,
polynomial in $\rstar^{-1}$.
\item $C_2(n,R,\rstar)$ in~\eqref{eq:fraenkel-centroid}: Fraenkel
estimate at the centroid.
\item $C_3(n,R,\rstar)$ in~\eqref{eq:normal-osc-centroid}: normal
oscillation.
\item $C_4(n,R,\rstar)=\max\{n/\rstar,3R/(2\rstar)\}$ in~\eqref{eq:radial-trace}:
\textbf{the oriented $L^1$ radial trace constant}, the load-bearing
output of Theorem~\ref{thm:radial-trace}.
\item $C_5(n,R,\rstar),C_6(n,R,\rstar)$ in Corollaries~\ref{cor:T-pointwise},
\ref{cor:H-trace-good}: levelwise inputs to \S4.4.
\end{itemize}
Every constant is polynomial in $n,R,\rstar^{-1}$ and independent of
$\delta_T(\Om)$ and of the boundary-layer constant $\kappa$.

%---- end of plan2_fj_center_radial_trace.tex ----
